\chapter*{How to read this book}
\addcontentsline{toc}{chapter}{How to read this book}
\markboth{How to read this book}{How to read this book}

\section*{What this is}

This book develops \emph{Cognitive Space Theory}: an account of structured
meaning in which a cognitive state is a \emph{section of a cellular sheaf}
over a simplicial complex, and in which the ways a cognitive state can fail
are \emph{cohomological} --- classified, localised and counted by the same
algebra that produces them.

The theory is under development. It is presented here as mathematics, not as
a specification: no programming language appears, no system is described, and
no result depends on a particular implementation. Where a construction is a
design intention rather than a theorem, it says so in the margin, in colour,
on the line where the claim is made.

\section*{Why every claim carries a tag}

A theory under development contains, at any moment, a mixture of proved
statements, statements proved only under hypotheses one may not be able to
meet, statements that have been computed but not proved, statements that are
intentions, and statements that were believed and are now known to be false.
Written in a single register these are indistinguishable, and a reader
studying such a text learns falsehoods with the same confidence as theorems.

So every substantive claim here carries one of these:

\begin{center}
\begin{tabular}{@{}>{\raggedright\arraybackslash}p{0.38\textwidth}%
                  >{\raggedright\arraybackslash}p{0.55\textwidth}@{}}
\toprule
\Proved & A complete proof is given here, from stated hypotheses.\\
\addlinespace[2pt]
\ProvedCond & Proved, but only under hypotheses that may not be available.\\
\addlinespace[2pt]
\Measured & A number produced by a procedure that can be re-run. Not a proof.\\
\addlinespace[2pt]
\Implemented & Computable as stated; no proof of correctness given.\\
\addlinespace[2pt]
\Proposed & A design intention. Not proved.\\
\addlinespace[2pt]
\Contested & Disputed on stated grounds, in place.\\
\addlinespace[2pt]
\Refuted & Shown false, or tested and failed. Retained \emph{with} the refutation.\\
\bottomrule
\end{tabular}
\end{center}

Refuted material is kept rather than deleted. Knowing why a construction
fails is most of what is needed in order not to rebuild it.

\section*{The three layers}

Each central object is presented three times over, so that a single reading
can be shallow or deep:

\begin{enumerate}
  \item a boxed \textbf{Intuition} in plain language, with no symbols;
  \item the formal \textbf{Definition}, \textbf{Construction} or
        \textbf{Proposition}, with proof;
  \item a \textbf{Worked example} on an explicit small complex, carried
        through to numbers.
\end{enumerate}

The worked examples are cumulative. A single running example --- a
four-vertex complex $\Kc_4$, which is \emph{derived} in
Chapter~\ref{ch:ingest} rather than postulated --- is reused throughout, so
that the same object is measured by every instrument the book builds. Where
an exercise appears, it is because the step is short and doing it yourself is
faster than reading it.

\section*{Agendas, and what an incomplete section looks like}

Every section of Parts~\ref{part:orientation} and~\ref{part:foundations}
ends with an \textbf{Agenda} box in green, containing three lists:

\begin{itemize}
  \item \textit{Settled} --- what the section establishes, by proof or by
        citation, and which may be used downstream without further argument;
  \item \textit{Open} --- what it does \emph{not} establish. These are the
        honest gaps, stated in the same place as the material that needs them;
  \item \textit{Next steps} --- the ordered work that would close the gaps.
\end{itemize}

An agenda whose \textit{Open} list is non-empty is a section that is
mathematically incomplete, and the incompleteness is printed rather than
implied. Later parts of the book, which have not yet been brought up to the
standard of Parts~\ref{part:orientation} and~\ref{part:foundations}, instead
carry a blue \textbf{Work plan} box. The two are not the same:

\begin{center}
\begin{tabular}{@{}>{\raggedright\arraybackslash}p{0.28\textwidth}%
                  >{\raggedright\arraybackslash}p{0.65\textwidth}@{}}
\toprule
\textsf{\textbf{\color{examplecol}Agenda}} (green) &
  Prose exists. It is correct as far as it goes, and the box says how far
  that is.\\
\addlinespace[3pt]
\textsf{\textbf{\color{linkblue}Work plan}} (blue) &
  Prose does \emph{not} exist, or predates the foundations. The box is the
  instruction for writing it.\\
\bottomrule
\end{tabular}
\end{center}

\section*{Reading order}

Part~\ref{part:orientation} is short, has no prerequisites, and states the
one hypothesis everything else is answerable to. Read it first.

Part~\ref{part:foundations} is the mathematical spine and is meant to be read
in order: the base complex is \emph{constructed} in
Chapter~\ref{ch:ingest} before it is used in Chapter~\ref{ch:complex}, the
sheaf is laid on it in Chapter~\ref{ch:sheaf}, and the operator that measures
everything appears in Chapter~\ref{ch:laplacian}. Chapter~\ref{ch:config}
then corrects a type error that the first four chapters deliberately carry:
the state is written as though the space it lives in were fixed, and it is
not.

A reader who wants the shortest honest path to the point of the theory should
read \S\ref{sec:hypothesis}, \S\ref{sec:degenerate} and
Theorem~\ref{thm:functoriality}. Those three say, respectively, why a complex
is necessary at all, what is lost by using a graph instead, and what the
higher cells buy in exact algebraic terms.
